\chapter{Adatbázis mélyebben: tranzakció, migráció és konkurencia}\label{ch:db-deep-dive-hu}

\noindent\textbf{Tanulási út: Gyors út: alap.} Tanuld meg a perzisztens alkalmazásokhoz szükséges séma- és query-invariánsokat.\par\medskip

\rustbridgeblockhu{A függvényszignatúrák, \code{Result}, \code{Vec<T>}, typed paraméterek és a \code{?} ismerősek maradnak.}{A \code{#[query]}, compiler-owned SQL, typed bind paraméter, statikus sorbound és \code{Db}/transaction capability teszi felülvizsgálhatóvá az adatbázis-effectet.}{Ne tanuld meg a stringből épített SQL-t Velran-technikaként. Az új fogalom a verified DB-határ, nem a queryk stringes összerakása.}


\invariantblock{Biztonsági invariáns: az adat nem SQL-szerkezet}{Az SQL szerkezete compiler-owned marad. Felhasználó által befolyásolt érték csak typed bind paraméterként kerülhet a querybe, az eredmény mérete pedig explicit bounded; a validáció nem teszi megengedetté a stringből épített SQL-t.}

\diagnosticblockhu{Compiler diagnostic}{SEC-A05-001}{A forrás unsafe SQL-összeállítást próbál használni compiler-owned SQL és typed bindok helyett.}{Az SQL szerkezete maradjon statikus; a requestből érkező értékeket typed bind paraméterként add át.}

Az előző fejezet bemutatta a typed SQL-t és a CRUD-ot. Itt azokra az adatbázis-szabályokra koncentrálunk, amelyek párhuzamos kérések és release-ek között védik az integritást. Első olvasáskor a query-szerződésre, typed bindokra, tranzakcióra, optimistic lockingra és migrationre figyelj; a kapacitás-, backup- és engine-választási részletek későbbi referencia.

\section{Három külön felelősség}

Velran alkalmazásnál érdemes három réteget szétválasztani:

\begin{description}
\item[alkalmazáskód] statikus typed SQL queryket deklarál, bindokat ad át és tranzakciót kér;
\item[runtime] kapcsolatot pooloz, bindol, timeoutot és eredménykorlátot alkalmaz, majd a visszakapott sorokat a deklarált típushoz ellenőrzi;
\item[operator] DB credentialt, TLS-t, adatbázis-kapacitást, backupot és migrációs folyamatot birtokol.
\end{description}

Az alkalmazás tehát nem nyit tetszőleges socketet az adatbázishoz és nem épít saját connection poolt. A \code{Db} és a \code{Transaction} capability az a felület, amelyen keresztül adatbázis-művelet történhet.

A V1 data layer általános defense-in-depth korlátai között van a legfeljebb 32 poolkapcsolat, az 5 másodperces connection-acquire timeout, a 10 másodperces query timeout, a 10\,000 soros backend eredményplafon és a 16 MiB-os eredményméret-plafon. A Velran modell-lista query ennél szigorúbb: a compiler literal \code{LIMIT}-et követel, és compiler/runtime szinten 1000 sor a felső korlát. Ezek biztonsági plafonok, nem pagination design; a helyes query ettől függetlenül legyen rövid, paginált és indexelhető.

Több szerverpéldánynál a DB connection budgetet az operatornak összesítve kell méreteznie:

\begin{lstlisting}
ossz DB kapcsolat ~= Velran instance-ok szama
                    * instance-onkenti pool
                    + admin/maintenance tartalek
\end{lstlisting}

\section{A query visszatérési típusa szerződés}

Egy lekérdezés cardinalityje már a függvény típusában látszik:

\begin{lstlisting}[language=Velran]
#[query]
fn articleById(db: Db, id: i64)
    -> Result<Article, DbError> sql {
    SELECT id, slug, title, body, version
    FROM articles
    WHERE id = :id
}

#[query]
fn articleBySlug(db: Db, slug: Slug)
    -> Result<Option<Article>, DbError> sql {
    SELECT id, slug, title, body, version
    FROM articles
    WHERE slug = :slug
}

#[query]
fn recentArticles(db: Db, offset: i64)
    -> Result<Vec<Article>, DbError> sql {
    SELECT id, slug, title, body, version
    FROM articles
    ORDER BY id DESC
    LIMIT 100 OFFSET :offset
}
\end{lstlisting}

A három forma jelentése:

\begin{center}
\begin{tabularx}{\textwidth}{@{}L{0.23\textwidth}YY@{}}
\toprule
Típus & Elvárt eredmény & Cardinality sérülése \\
\midrule
\code{Article} & pontosan egy sor & 0 sor: \code{404}; több sor: DB-hiba \\
\code{Option<Article>} & nulla vagy egy sor & több sor: DB-hiba \\
\code{Vec<Article>} & nulla vagy több sor & a runtime eredménylimitei érvényesek \\
\bottomrule
\end{tabularx}
\end{center}

Ez több egyszerű kényelmi jelölésnél: a lekérdezés eredményének alakja része a program futás közben is kikényszerített szerződésének. Az \code{Option<Article>} nem azt jelenti, hogy ``vedd az első sort, ha van'': ha a query kettő vagy több sort ad, az invariáns sérült, ezért a runtime fail-closed DB-hibát ad. Ugyanezért egy \code{Article} query nulla sor esetén \code{404 Not Found}, több sor esetén pedig nem választ önkényesen rekordot.

A \code{Vec<Model>} querynek az SQL-ben statikus literal sorlimitje is kell legyen. A bound a reusable query contract része, nem a hívó által adott \code{limit} paraméter; használj fix \code{LIMIT}-et, és bounded offsettel vagy más explicit hozzáférési mintával paginálj.

\subsection*{A row shape legyen explicit}

A \code{SELECT} mezőinek neve és sorrendje egyezzen a modellel. Ne használj \code{SELECT *}-ot alkalmazásqueryben.

\begin{lstlisting}[language=Velran]
model Article {
    id: i64
    slug: Slug
    title: String
    body: String
    version: i64
}

// jo: ugyanaz az ot mezo, ugyanebben a sorrendben
SELECT id, slug, title, body, version
FROM articles
\end{lstlisting}

Az explicit projection két okból fontos. Először is a fordító ellenőrizni tudja a modell és a query kapcsolatát. Másodszor egy későbbi schema-bővítés nem változtatja meg észrevétlenül egy régi query result shape-jét.

\section{Bindolj, ne SQL-t építs}

A kliensből származó érték typed bindként kerüljön az SQL-be:

\begin{lstlisting}[language=Velran]
#[query]
fn findByEmail(db: Db, email: String)
    -> Result<Option<User>, DbError> sql {
    SELECT id, email, display_name
    FROM users
    WHERE email = :email
}
\end{lstlisting}

A \code{:email} nem string-helyettesítés. A runtime backend-specifikus placeholderre fordítja és paraméterként bindolja.

Amit ne próbálj meg:

\begin{lstlisting}
// rossz gondolkodasmod:
// "SELECT ... ORDER BY " + userOrder
// "WHERE " + userFilter
\end{lstlisting}

Oszlopnév, \code{ORDER BY} irány, SQL operator vagy teljes WHERE-fragment ne érkezzen user inputból. Ha több rendezést akarsz támogatni, az alkalmazáskódban válassz néhány statikus query közül.

\section{Olvasás és mutáció authority-határa}

Olvasó query \code{Db} capabilityt kap. Adatmódosító query \code{Transaction}-t:

\begin{lstlisting}[language=Velran]
#[query]
fn renameArticle(
    tx: Transaction,
    id: i64,
    title: String
) -> Result<(), DbError> mutates Article by id sql {
    UPDATE articles
    SET title = :title
    WHERE id = :id
}
\end{lstlisting}

A hívó action explicit tranzakciós blokkot nyit:

\begin{lstlisting}[language=Velran]
#[action]
fn rename(
    ctx: ActionContext,
    db: Db,
    id: i64,
    title: String
) -> Result<Redirect, PageError> {
    let article = articleById(db, id)?;
    authorize article authenticated;
    transaction db {
        renameArticle(tx, article.id, title)?;
    }
    return Ok(redirect(adminArticles()));
}
\end{lstlisting}

Normál best-effort alkalmazáskódban a sikeres blokk commitol, ismert DB hiba rollbackel. Retry-szenzitív vagy critical műveletnél viszont ne egyszerű siker/hiba állapotra szűkítsd az eredményt. Capture-öld explicit a transaction outcome-ot:

\begin{lstlisting}[language=Velran]
let article = articleById(db, id)?;
authorize article authenticated;
let outcome = transaction db {
    renameArticle(tx, article.id, title)?;
};

match outcome {
    Committed => { return Ok(redirect(adminArticles())); }
    RolledBack => { fail conflict; }
    CommitUnknown => { fail conflict; }
}
\end{lstlisting}

A \code{CommitUnknown} azt jelenti, hogy a runtime nem kapott elég bizonyítékot annak eldöntésére, megtörtént-e a commit. Ez nem válhat automatikus retry-já. Critical idempotens transactionnél a compiler kötelezővé teszi az outcome capture-t és exhaustive match-et. A \code{Transaction} továbbra sem escape-elhet a blokkból.

\subsection*{Egy üzleti művelet legyen egy tranzakció}

Több, összetartozó módosítást ne külön actionök vagy külön tranzakciók végezzenek:

\begin{lstlisting}[language=Velran]
transaction db {
    Article.publishChanged(tx, article.id, version)?;
    audit Article article.id action publish;
        from article.status to ArticleStatus.Published
}
\end{lstlisting}

Ebben a mintában a cikk publikálása és a business audit trail együtt commitol. Ha az audit nem írható, a publikálás is rollbackel. Így nem jön létre olyan állapot, amelyben az üzleti módosítás megtörtént, de a hozzá kötelező audit rekord hiányzik.

\section{Optimistic locking: ne írd felül más munkáját}

CMS-ben és wikiben gyakori helyzet, hogy két szerkesztő ugyanazt a rekordot nyitja meg. A klasszikus

\begin{lstlisting}
UPDATE articles SET title = :title WHERE id = :id
\end{lstlisting}

az utoljára mentő felhasználót teszi nyertessé, és csendben elveszhet az első szerkesztő módosítása.

Ehhez használható verziómező és \code{Changed}:

\begin{lstlisting}[language=Velran]
model Article {
    id: i64
    title: String
    version: i64
}

#[query]
fn updateArticle(
    tx: Transaction,
    id: i64,
    title: String,
    version: i64
) -> Result<Changed, DbError> mutates Article by id sql {
    UPDATE articles
    SET title = :title,
        version = version + 1
    WHERE id = :id AND version = :version
}
\end{lstlisting}

A call site előbb töltse be és autorizálja a cikket, és a request ID helyett az \code{article.id} értéket adja át ennek a mutációnak. Az edit form elküldi azt a \code{version} értéket, amelyet a felhasználó az oldal megnyitásakor látott. A \code{Changed} szerződése szigorú: a mutációnak \textbf{pontosan egy sort} kell módosítania. Egy módosított sor siker; nulla módosított sor \code{409 Conflict}; kettő vagy több módosított sor DB-hiba, mert a query a saját invariánsát sértette meg. \code{Changed} csak mutáló query visszatérési típusa lehet, és V1-ben nem használható \code{RETURNING} klauzussal.

Optimistic lockingnál, ha időközben más már mentett, a \code{WHERE id = :id AND version = :version} nulla sort talál, ezért a kérés konfliktusra fut. A helyes UX ilyenkor nem az automatikus újrapróbálás. Mutasd meg, hogy az adat közben változott, töltsd be az új verziót, és a felhasználó döntsön az összeolvasztásról.

Optimistic locking tipikus helyei:

\begin{itemize}
\item wiki oldalak és CMS cikkek;
\item ticket vagy ügyfélrekord;
\item workflow-állapotot tartalmazó üzleti objektum;
\item ritkán módosított admin konfiguráció.
\end{itemize}

Append-only eseményekhez vagy adatbázis-oldali atomi számlálókhoz általában nincs rá szükség.

\section{Indexet a hozzáférési mintához tervezz}

A jó index nem attól jó, hogy sok van belőle, hanem attól, hogy a tényleges queryt szolgálja ki.

Egy CMS-ben például természetes hozzáférési minták lehetnek:

\begin{lstlisting}[language=SQL]
CREATE UNIQUE INDEX articles_slug_uq
    ON articles(slug);

CREATE INDEX articles_status_id_idx
    ON articles(status, id);
\end{lstlisting}

Ha a publikus oldal slug alapján kér le egy cikket, a slug legyen unique/indexelt. Ha az admin lista státusz szerint szűr és azon belül azonosító szerint lapoz, ehhez illeszkedő összetett index kellhet.

Néhány szabály:

\begin{itemize}
\item minden nagy listát paginálj;
\item a rendszeresen használt lookup mező legyen indexelhető;
\item az indexek sorrendje kövesse a WHERE/ORDER BY mintát;
\item ne indexelj automatikusan minden oszlopot: az index írási és storage költség;
\item teljesítményproblémát valós backend \code{EXPLAIN} tervével vizsgálj.
\end{itemize}

A Velran typed SQL-ja a query alakját védi; az adatbázis optimalizálóját nem helyettesíti.

\section{Migration: a séma a release része}

Productionban a server nem futtat automatikus migrációt startupkor. Ez szándékos: egy schema-változás külön operátori művelet, külön credentialdel és külön hibahatárral.

Tipikus könyvtár:

\begin{lstlisting}
migrations/
  0001_create_articles.sql
  0002_articles_slug_index.sql
  0003_articles_add_version.sql
\end{lstlisting}

A fájlnév alakja:

\begin{lstlisting}
NNNN_name.sql
\end{lstlisting}

Már alkalmazott migration fájlt ne módosíts. A CLI SHA-256 checksumot tárol, és eltérésnél megáll.

A production parancsok a telepített CLI-val:

\begin{lstlisting}
/usr/local/bin/velran-cli migrate status \
  --dir /srv/myapp/current/migrations \
  --db-url-file /run/secrets/migration-db-url

/usr/local/bin/velran-cli migrate verify \
  --dir /srv/myapp/current/migrations \
  --db-url-file /run/secrets/migration-db-url

/usr/local/bin/velran-cli migrate apply \
  --dir /srv/myapp/current/migrations \
  --db-url-file /run/secrets/migration-db-url \
  --lock-timeout-secs 30
\end{lstlisting}

A migration credential magasabb jogosultságú lehet, mint az alkalmazás DML credentialje. Emiatt a server process ne kapja meg.

\subsection*{Mit garantál a migration CLI?}

\begin{itemize}
\item ellenőrzi az alkalmazott fájl nevét és checksumát;
\item nem enged észrevétlenül régi verziószám alá új migrationt beszúrni;
\item adatbázis-szintű migration lockot kér;
\item a state csak sikeres alkalmazás után halad tovább.
\end{itemize}

Backend-specifikus különbség viszont marad. PostgreSQL migrationonként tranzakciót használ; SQLite \code{BEGIN IMMEDIATE} writer lockkal a batch-et tranzakcióban tartja; MariaDB DDL nem minden esetben tranzakcionális, ezért részleges DDL-hibánál operator repair is szükséges lehet.

\section{Expand/contract: tartsd meg a rollback lehetőségét}

Egy gyors application rollback csak akkor egyszerű, ha az új schema a régi programmal is kompatibilis.

Ezért nagyobb változtatásnál használj expand/contract mintát:

\begin{enumerate}
\item \textbf{expand}: adj hozzá új nullable/defaultolt mezőt vagy új táblát úgy, hogy a régi app még működjön;
\item deployold az új alkalmazást;
\item szükség esetén végezz kontrollált backfillt;
\item csak későbbi release-ben \textbf{contract}: töröld a régi mezőt vagy constraintet, amikor a rollback window lezárult.
\end{enumerate}

Az azonnali rename/drop, új kötelező mező vagy inkompatibilis adattípus-váltás könnyen forward-only release-t eredményez. Ezt release review során explicit döntésként kezeld.

\section{Adatbázis-hiba kezelése}

A \code{DbError} technikai hiba. A végfelhasználónak ne jeleníts meg driver message-et, SQL szöveget, hostnevet vagy connection stringet.

Alkalmazási szinten különítsd el legalább ezeket a helyzeteket:

\begin{description}
\item[nincs ilyen rekord] normál domain-állapot, opcionális queryvel kezelhető;
\item[409 konkurens módosítás] optimistic-locking konfliktus, a felhasználó újratöltéssel/összevetéssel oldja fel;
\item[adatbázis nem elérhető] infrastruktúrahiba; readiness is jelezheti, a logban request ID-val kereshető;
\item[constraint/domain hiba] inkább validációval előzd meg, de az adatbázis constraint maradjon utolsó integritási védővonal.
\end{description}

Ne építs automatikus, korlátlan retry loopot egy teljes mutáló action köré. Egy megismételt kérésnek mellékhatása lehet, és a tranzakciós konfliktus sem feltétlen átmeneti hiba.

\section{SQLite, PostgreSQL vagy MariaDB?}

Velran V1 mindhárom backendet támogatja, de ez nem jelenti azt, hogy minden SQL dialektusfeature hordozható.

\begin{description}
\item[SQLite] kiváló egyszerű vagy egygépes telepítéshez; az application DB fájl backupját SQLite-aware módszerrel készítsd, ne élő fájl egyszerű másolásával.
\item[PostgreSQL] erős alap többpéldányos és nagyobb production rendszerekhez; TLS-t használj remote kapcsolathoz.
\item[MariaDB] szintén támogatott, de különösen migrációknál számolj a DDL tranzakciós különbségeivel.
\end{description}

Portable alkalmazásnál a queryk közös SQL-részhalmazra épüljenek. Backend-specifikus funkciót csak tudatosan használj, és akkor a projekt dokumentációjában rögzítsd, hogy az alkalmazás mely backendhez kötött.

\section{Backup szempont adatbázis-fejlesztőként}

Adatbázis-fejlesztőként itt egy dolgot kell rögzíteni: a migráció és a backup két külön életciklus. A migration leírja, hogyan változik a séma; a backup/restore azt bizonyítja, hogy a tartós állapot visszaállítható. PostgreSQLnél, MariaDB-nél és SQLite-nál a backend saját konzisztens mentési eszközét használd, és ne tekints egy egyszerű fájlmásolást automatikusan érvényes backupnak.

A teljes alkalmazásállapot -- adatbázis, AppFs, local-auth store, recovery manifest, recovery point objective (RPO), recovery time objective (RTO) és restore drill -- a \ref{ch:recovery}. fejezet témája. Itt azért utalunk rá, mert egy schema-változtatás tervezésekor már tudni kell, hogy rollbackhez vagy restore-hoz milyen állapot áll rendelkezésre.

\section{Egy CMS adatbázis-rétegének mintája}

Egy kisebb CMS-nél jó kiindulás:

\begin{lstlisting}
Article
  id
  slug          UNIQUE
  title
  body
  status
  hero_image
  version
  created_at
  updated_at

migrations/
  0001_create_articles.sql
  0002_article_indexes.sql
  0003_add_version.sql
\end{lstlisting}

Az alkalmazási minták:

\begin{itemize}
\item publikus olvasás: \code{Option<Article>} slug alapján;
\item admin lista: paginált \code{Vec<Article>};
\item szerkesztés: \code{Changed} + \code{version};
\item publish/archive: egy transactionben üzleti módosítás + audit;
\item schema-változás: külön CLI migration;
\item deploy előtt: migration verify + backup/restore readiness;
\item deploy után: readiness, log, metrics és smoke test.
\end{itemize}

Ezzel a modellel a database layer nem különálló SQL-gyűjtemény lesz, hanem a webalkalmazás típus-, konkurencia-, release- és recovery-szerződésének része.

\section{Fejezet végi ellenőrzőlista}

Mielőtt egy adatbázisos funkciót késznek tekintesz, ellenőrizd:

\begin{itemize}
\item minden user érték typed bindként kerül SQL-be;
\item nincs \code{SELECT *}; a projection egyezik a modellel;
\item nagy lista paginált;
\item a fontos lookup és rendezési minta mögött van megfelelő index;
\item minden mutáció \code{Transaction} capabilityvel fut;
\item több összetartozó módosítás egy tranzakcióban van;
\item ember által párhuzamosan szerkeszthető rekordnál van optimistic locking;
\item már alkalmazott migration immutable;
\item a migration credential nincs az application servernél;
\item destructive schema-változásnál ismert a rollback következménye;
\item a DB hiba nem szivárogtat SQL-t vagy credentialt a kliensnek;
\item a backuphoz tartozik rendszeresen bizonyított restore.
\end{itemize}

\section{Gyakorlat: kezeld a query-t interfészként}
\paragraph{Gyakorlási mód: támogatott.} Használd az előszóban meghatározott támogatott módszert.

Minden feature query-jénél írd le az inputtípusokat, a projection oszlopait, a várt kardinalitást, a rendezést és az index-feltételezést. Ezt interfész-contractként kezeld, ne SQL-be rejtett implementációs részletként.

\section{Tipikus hiba: az adatbázist csak tárolónak tekintjük}
A production adatbázis invariánsokat és concurrency-szemantikát is kikényszerít. Ha a helyesség attól függ, hogy egy feltétel párhuzamos kérések mellett is igaz maradjon, az authoritative részt az adatbázisban vagy a tranzakciós tervben rögzítsd, ne csak egy korábbi alkalmazásszintű ellenőrzésben.

\section{Folyamatos mintaprojekt: Secure Notes}
Indulj az \path{examples/secure-notes/migrations/0001_init.sql} sémából. Vizsgáld felül a cím/törzs korlátokat, a publikációs állapot constraintjét, az időbélyegeket és az owner/id indexet mint adatbázisszintű invariantokat. Ezután adj lapozást úgy, hogy a determinisztikus rendezés és a principal szerinti szűrés megmaradjon; dokumentáld, mely garancia tartozik az SQL/séma réteghez, és melyik marad handler-level policy.
